% !TEX TS-program = pdflatex
% !TeX program = pdflatex
% !TEX encoding = UTF-8
% !TeX spellcheck = en_US

\documentclass[11pt, a4paper]{article}
%\usepackage{fullpage}
\usepackage[left=1cm,right=1cm,top=1cm,bottom=2cm]{geometry}
\usepackage[fleqn]{amsmath}
\usepackage{amssymb}
%\usepackage{indentfirst}
\usepackage[T1]{fontenc}
\usepackage[utf8]{inputenc}
\usepackage[french,english]{babel}
\usepackage{txfonts} 
\usepackage[]{graphicx}
\usepackage{multirow}
\usepackage{hyperref}
\usepackage{parskip}
\usepackage{multicol}
\usepackage{wrapfig}

\usepackage{turnstile}%Induction symbole

\usepackage{tikz}
\usetikzlibrary{arrows, automata}
\usetikzlibrary{decorations.pathmorphing}

\renewcommand{\baselinestretch}{1}

\setlength{\parindent}{24pt}


\begin{document}

%\selectlanguage {french}
%\pagestyle{empty} 

\noindent
\begin{tabular}{ll}
\multirow{3}{*}{\includegraphics[width=1.5cm]{../../../extra/logo/esi.nlp.pdf}} & 
\'Ecole national Supérieure d'Informatique\\
& 2\textsuperscript{nd} year second cycle (2CSSID)\\
& NLP: Natural Language Processing (2022-2023)
\end{tabular}\\[.25cm]
\noindent\rule{\textwidth}{2pt}\\[-0.5cm]
\begin{center}
{\LARGE \textbf{Lab01: Naive Spell Checker}}
\begin{flushright}
	Abdelkrime Aries
\end{flushright}
\end{center}\vspace{-0.5cm}
\noindent\rule{\textwidth}{2pt}

We want to implement a spell checker from scratch. 
In our first attempt, we use a dictionary to look up the word.
If it is found, then there is no need to spell-check it.
Otherwise, we look for similar words in the said dictionary.
To test for similarity, we use Levenstein distance; the least distant words are the most similar ones.
The following section better explain this process.


\section*{1. Program description}

Here, different functions are described; either those implemented or not. 
You have to understand how they are implemented to respond on the different questions.

\subsection*{1.1. Extracting words}

This function extracts a list of words from a dictionary given a word.
Its extracts words which have the same first letter or second letter as the given word.

\textbf{\slshape This is already implemented}.

\subsection*{1.2. Levenstein measure}

This function calculates levenstein distance between two given words.
In addition to the distance, it returns the dynamic table.
If one of the words is empty, it returns the length of the other word as levenstein and an empty dynamic table.

\textbf{\slshape This must be implemented}.

\subsection*{1.3. Choosing similar words}

This function receives a list of possible words, a target word and a threshold.
Among these words it must select the most similar ones based on a similarity measure and a threshold.
The similarity measure between two words $w_1$ and $w_2$ is given as follows :
\[sim(w_1, w_2) = \frac{\max(|w_1|, |w_2|) - levenstein(w_1, w_2)}{\max(|w_1|, |w_2|)}\]

\textbf{\slshape This must be implemented}.

\subsection*{1.4. Checking a word}

Given a word, 
\begin{itemize}
	\item This function will call for the function responsible for extracting possible words from a list.
	\item Then, it checks if the word exists in this list. In this case it returns a list containing the same word.
	\item Otherwise, it passes this list into the function responsible for choosing similar words with a threshold of 75\% and a substitution cost of 1.
	\item It returns the result.
\end{itemize}

\textbf{\slshape This is already implemented}.

\section*{3. Questions}

Answer these questions at the end of your code, as comments.
\begin{enumerate}
	\item We noticed that some words are not taken in consideration although they have more chance.
	For example, the word \textbf{qsministrator} against \textbf{administrator}. 
	Why? Propose an amelioration.
	
	\item Some words are not detected as correct although they are.
	Example, the word \textbf{confirmation} is correct, but detected as a spell error.
	The dictionary contains only the word \textbf{confirm}.
	Propose a solution without changing the content of the dictionary.
	
	\item We want to evaluate our system. 
	For each word, we prepared a list of possible corrections ordered by order of similarity.
	Propose a method to evaluate such system.
	
\end{enumerate}


\section*{4. Evaluation}

\begin{itemize}
	\item Duration: 1h (You have to return your work at the end of session)
	\item Grade
	\begin{itemize}
		\item \textbf{levenstein grade} (8pts) = zero length (1pt) + initialize first row (1pt) + initialize first column (1pt) + character deletion (1pt) + character adding (1pt) + character substitution (2pts) + correctness test (1pt).
		\item \textbf{choose words grade} (7pts) = selecting top words (2pts) + sorting them (2pts) + splitting them into words and scores (2pts) + in case there are no suggestions (1pt).
		\item \textbf{questions grade} (3pts) = 1pt for each question.
		\item \textbf{In time grade} (2pts)
	\end{itemize}
\end{itemize}

\end{document}
